\section{External reference qualification overview}
Strict independent-reference accuracy remains unresolved. Table~\ref{tab:external-qualification} keeps complete or bounded reference years, stations meeting each history rule, quality and clock, independence, and scored cases in separate fields. Nine Greater Wellington sites stayed at the metadata screen because temperature series were not requested, so their complete hot-year count is unknown \cite{greaterwellington2016quality}. The saved screens found zero strict complete years at 24 NEON sites and zero complete July--June thermal years at five R\'eunion sites; eight new USGS candidates contain \USGSCompleteYears{} complete station-years, but \USGSFourYearStations{} stations reach four years \cite{marzolf2025neon,hubeautemp2026,usgsstatistics2026}. That USGS count is not an hourly physical reference and does not prove independence. ECCC Red River shares a river with a source station and retains grade and time-basis questions \cite{ecccautomated2024,usgsred2025}. An Australian candidate archive had prior evaluation exposure: \AUFirstEvaluable{} of \AURefYears{} candidate station-years were previously evaluable, including all \AUFormerTestEvaluable{} designated test years \cite{bom2026}. These exclusions do not establish that the observations are erroneous.

\begin{table}[p]
\centering
\scriptsize
\setlength{\tabcolsep}{3pt}
\renewcommand{\arraystretch}{1.08}
\caption{External reference qualification. Complete or bounded reference years, stations meeting the stated history rule, quality and clock, independence, and scored cases are separate fields. Unknown is written as unknown. A zero appears only when the saved screen counted that quantity.}
\label{tab:external-qualification}
\begin{tabular}{@{}p{0.13\textwidth}p{0.16\textwidth}p{0.16\textwidth}p{0.16\textwidth}p{0.15\textwidth}p{0.16\textwidth}@{}}
\toprule
Source & Reference years & History-rule stations & Quality and clock & Independence & Scored cases \\
\midrule
Greater Wellington & Temperature series not acquired & 9 potential sites; quality-code screen not completed & Quality-code values were not requested & Unknown & Not scored \\
NEON & 0 strict complete years & 0 of 24 sites meet the four-year strict rule & UTC day and final quality flag required & Climate independence not established & Not scored \\
New USGS candidates & 8 complete station-years & 0 of 8 sites reach four complete years & Daily screen only; hourly reference not verified here & Not established & Not scored. The observation-only witness check is separate. \\
R\'eunion (Hub'Eau) & 0 complete July--June thermal years & 0 of 5 sites meet the four-year rule & Reported-day time zone unresolved & Unknown & Not scored \\
ECCC Red River & 0 complete local wall-clock years & Strict physical reference not met & Grade and time basis unresolved & Shares a river with a source station & Not scored \\
Australia candidate & 80 candidate station-years; strict 24-hour sampling not established & 79 of 80 were previously evaluable, including 10 of 10 designated test years & The saved screen allowed 20 observed hours & Not independent: prior evaluation exposure & Not scored as a new test \\
British Columbia & 16 published station-years; 0 strict 24-hour stations & 4 stations scored; 0 meet the strict 24-hour rule & Days may contain 20 hours; clock unresolved & Not established; three stations share the Fraser basin & 372 masks total: 288 controlled gaps and 84 calendar gaps. The reported MAEs use the 288. \\
Nova Scotia & 8 reported-clock station-years & 2 stations; strict reference not met & Source QC and time zone unresolved & Not established & 144 controlled-gap masks scored \\
Truckee, California & 4 reported years & 1 station; strict reference not met & Clock and estimation flags unresolved & Not established; one managed station with prior basin contact & 72 controlled-gap masks scored \\
USGS Rio Mora & Complete quarter-hour year: 2017. Bounded years: 2012 and 2019 (2012: 3 unknown slots; 2017: 0 unknown slots; 2019: 2 unknown slots) & 1 station & Unknown quarter-hour slots are assigned to the conditional range $[-5,50]\,^{\circ}$C & Not an independent climate & 54 hourly-supported masks scored against the reported daily target and the conditional continuous target \\
Reservoir network & 0 complete finite years among 69 station-years & 0 of 10 files have four finite years & Strict quarter-hour calendars were not met. The scored bounds assume missing values lie in $[-5,50]\,^{\circ}$C & Not established & 306 exploratory bounded-reference cases scored at 4 stations and 17 target station-years \\
\bottomrule
\end{tabular}
\end{table}


Two exploratory comparisons bound what external evidence can show. British Columbia has \BCStations{} stations and \BCYears{} station-years. Of \BCMaskTotal{} masks, \BCCases{} are controlled gaps and \BCCalendarMasks{} are calendar gaps. Plug-in, D2 and source-role-calibrated stochastic H30 mean absolute errors on the controlled gaps are \BCPluginMAE{}, \BCDTwoMAE{} and \BCStochasticMAE{} $^{\circ}$C \cite{moore2024bc,callahan2025air2stream}. Three stations share the Fraser basin; daily means may use only 20 hourly observations, their day convention is unresolved, and source quality processing used water--air comparisons \cite{wsc_temperature_faq}. Nova Scotia covers \NSExpStations{} stations, \NSExpYears{} station-years and \NSExpCases{} masks selected after inspecting availability; D2 and stochastic mean absolute errors are \NSExpDTwoMAE{} and \NSExpClassicMAE{} $^{\circ}$C \cite{novascotiawater2026}. Neither comparison is strict independent-climate validation.

The continuous USGS support at Rio Mora is the clearest warning about target definition. Three years have at least one approved unqualified continuous observation in every fixed-standard-time hour. Only \RioCompleteQuarterYear{} has every quarter-hour slot. The continuous targets for \RioBoundedYears{} assign unknown quarter-hour values to the conditional range $[-5,50]\,^{\circ}$C; they are not fully verified physical years. The \RioHourlyMasks{} hourly-supported masks give D2 MAE 0.180$^{\circ}$C against the published daily target. Using the original histories, D2 interval coverage is 53/54 against published daily H30 and 33/54 against that conditional continuous target; qualifier-aware histories yield 48/54 and 28/54. The difference can combine clock, rounding, corrections, daily preprocessing and the bound assumption. It does not identify a single source error or justify a universal physical calibration. Truckee scores \TruckeeCases{} controlled-gap masks on one station and four reported years. The reservoir comparison scores \ResScoreCases{} exploratory bounded-reference cases at \ResScoreStations{} stations and \ResScoreYears{} target station-years. Both remain descriptive checks of bounds and clock assumptions, not independent accuracy tests.

